% 論文節の日本語訳（読み物 PDF）用 preamble
% lualatex + jlreq
\documentclass[12pt,paper=a4]{jlreq}
\usepackage{luatexja-fontspec}
\usepackage{amsmath,amssymb,bm}
\usepackage{graphicx}
\usepackage{hyperref}
\usepackage{enumitem}
\usepackage{booktabs}
\usepackage{microtype}
\usepackage{xcolor}
\usepackage{caption}
\usepackage{tcolorbox}
\tcbuselibrary{breakable,skins}

\hypersetup{
  colorlinks=true,
  linkcolor=blue!60!black,
  urlcolor=blue!60!black,
  citecolor=blue!60!black,
}

\graphicspath{{images/}{./}}

\setlength{\parindent}{1em}
\setlength{\parskip}{0.55em}
\setlist[itemize]{leftmargin=1.6em,itemsep=0.25em,topsep=0.35em}
\setlist[enumerate]{leftmargin=1.8em,itemsep=0.25em,topsep=0.35em}

\captionsetup{
  font=small,
  labelfont=bf,
  skip=0.4em,
  justification=raggedright,
  singlelinecheck=false,
}

\newcommand{\vx}{\bm{x}}
\newcommand{\vu}{\bm{u}}
\newcommand{\veps}{\bm{\epsilon}}
\newcommand{\vw}{\bm{w}}

\DeclareMathOperator*{\argmax}{arg\,max}
\DeclareMathOperator*{\argmin}{arg\,min}
\DeclareMathOperator{\Tr}{Tr}
\DeclareMathOperator{\diag}{diag}

% 命題・定理ボックス（ページ跨ぎ時も切れ目を抑える）
\newtcolorbox{propositionbox}[1][]{%
  enhanced,
  breakable,
  colback=blue!4!white,
  colframe=blue!45!black,
  coltitle=white,
  fonttitle=\bfseries\small,
  title={#1},
  attach boxed title to top left={yshift=-2mm, xshift=3mm},
  boxed title style={
    colback=blue!45!black,
    colframe=blue!45!black,
    boxrule=0pt,
    arc=1pt,
  },
  boxrule=0.5pt,
  arc=1.5pt,
  left=5pt, right=5pt, top=8pt, bottom=5pt,
  before skip=0.9em,
  after skip=0.9em,
  pad at break=3mm,
}

% 図はフロートせず本文位置に置く（分断を防ぐ）
\newcommand{\translationfigure}[2]{%
  \begin{center}
    \includegraphics[width=0.88\linewidth]{#1}\\[0.35em]
    {\small #2\par}
  \end{center}
  \vspace{0.4em}%
}

\begin{document}
{\color{gray}\footnotesize 出典: §7.1\par}
\vspace{0.2em}

\section*{7.1 分布間変換の序論}


拡散モデル（diffusion models）は終端分布を標準ガウス分布 $p_{\mathrm{prior}}$ に固定する。しかし多くの応用では、分布間変換（distribution-to-distribution translation）——すなわち、ソース分布 $p_{\mathrm{src}}$ を別のターゲット分布 $p_{\mathrm{tgt}}$ へ変換すること——が必要である。例としては、スケッチを写実的な画像へ変換することや、芸術様式間の変換が挙げられる。

現代の拡散手法はこの目的のための実用的な手段を提供する。SDEdit（Meng et al., 2021b）のような一端点手法（one-endpoint methods）は、$t=0$ のソース画像から始め、それを中間ステップ $t$ まで拡散させ、その後ターゲット領域向けに事前学習された拡散モデルを用いて過程を逆転させる。これにより、ターゲット分布の様式と内容に一致する出力が得られる。

Dual Diffusion Bridge（Su et al., 2022）のような両端点手法（two-endpoint methods）は、代わりに、典型的には $t=1$ におけるガウス分布である共有潜在分布を通じて両領域を接続する。前方確率流 ODE（forward probability-flow ODE）が $p_{\mathrm{src}}$ からこの潜在空間へサンプルを輸送し、ターゲット領域上で学習された逆 ODE がそれらを $p_{\mathrm{tgt}}$ へ写し戻す。こうしたサンプリング時アプローチを超えて、第5.2節で述べた Flow Matching の枠組みは、訓練に基づく代替を提供する。すなわち、$p_{\mathrm{src}}$ から $p_{\mathrm{tgt}}$ へ質量を連続的に動かす ODE 流を直接学習する。

重要な点として、分布間の変換には、別々に訓練された二つのモデル以上のものが必要である。両端点での力学を整合させ、しかも「最も安い」（費用効率の良い）やり方でそれを行う、原理的な写像が求められる。

本節では、拡散に基づく翻訳応用の数多くを概観するのではなく、この古典的な分布間問題の数学的基礎へ焦点を移す。とりわけ最適輸送（optimal transport, OT）とそのエントロピー変種であるシュレーディンガー橋（Schrödinger bridge, SB）を強調する。これらは、理論コミュニティにおいて、費用効率の良い（数学的な意味で）分布変換の標準的定式化として長く研究されてきた。

その核心にある根本的な問いは次である。

\textbf{問い 7.1.1}

\begin{propositionbox}[注意]
二つの確率分布が与えられたとき、総費用を最小化しつつ一方を他方へ変換する最も効率的な方法は何か？
\end{propositionbox}


ここで費用 $c(x,y)$ は、点 $x$ から点 $y$ へ質量の単位を動かすことに対する罰則を割り当てる非負関数である。よくある選択は二乗距離 $c(x,y)=\|x-y\|^2$ である。

本節は、フローマッチングを含む拡散に基づくアプローチが、古典的および正則化された最適輸送とどのように結び付くかを明確にするための短い概観を与える。探求したい中心的な問いは次である。

\textbf{問い 7.1.2}

\begin{propositionbox}[注意]
拡散モデルは $p_{\mathrm{data}}$ と $p_{\mathrm{prior}}$ をつなぐ最適輸送の一形態であり、それはどのような意味においてか？
\end{propositionbox}


この問いに取り組むため、まず第7.2節で「最適性」が何を意味するかを明確にする。静的な Monge–Kantorovich 形式（式 (7.2.1)）における古典的最適輸送（OT）と、その動的な Benamou–Brenier 定式化（式 (7.2.3)；連続の方程式の下で運動エネルギーを最小化するもの）、ならびにエントロピー正則化変種（entropic OT）を式 (7.2.5) で復習する。後者はシュレーディンガー橋問題（式 (7.2.8)）と等価である。動的な見方では、OT は連続の方程式を満たす決定論的な流れを誘導するのに対し、SB は周辺分布が Fokker–Planck 方程式に従って発展する制御拡散を誘導する。これらの定式化間の高水準の対応を第7.3節で与える。

その後、議論を二部に分ける。まず第7.4節では、標準的な拡散モデルで用いられる固定の前方ノイズ付加 SDE それ自体は、任意の $p_{\mathrm{src}}$ と $p_{\mathrm{tgt}}$ の間のシュレーディンガー橋ではないことを説明する。前方過程は選ばれた参照拡散であり、時間前進および時間逆の両 SDE は一般に、所定のターゲットへの厳密な端点整合を強制しない。したがって、それらの端点について SB 問題を明示的に解かない限り、それはエントロピー OT 最適ではない。一方で、一方の出発点に固定されているという意味で、半橋（half-bridge）問題に対する最適解ではある。

次に第7.5節では、$p_{\mathrm{src}}=p_{\mathrm{prior}}$（ガウス）かつ $p_{\mathrm{tgt}}=p_{\mathrm{data}}$ という生成設定に戻る。PF-ODE は構成上 $p_{\mathrm{prior}}$ を $p_{\mathrm{data}}$ へ輸送する決定論的写像を定める。しかしこの流れは一般に、所定の輸送費用（例：二次 $W_2$）に対する OT 写像ではない。それは多数の許容可能な決定論的カップリングのうちの一つを実現するにすぎず、Benamou–Brenier 作用を最小化しない。続く議論では、「整流流（rectify flow）」手続き（第5.4.1節）が OT 写像に至るかどうかを検討するが、一般にはそのような理論的保証はない。したがって、拡散モデルの PF-ODE 写像と OT の厳密な特徴付けは、依然として挑戦的で未解決の問題である。
\end{document}
